\clearpage
\SourceLeaf{83}
\Needspace{10\baselineskip}\begin{center}
\textbf{NEUER BEWEIS DES LAGRANGISCHEN LEHRSATZES.}
\end{center}

\noindent\hspace*{\parindent}Der wichtige Lehrsatz des Herrn \textsc{de la Grange}, wonach jede Function von
$x$ durch eine Reihe nach den Potenzen von $u$ dargestellt wird, wenn $x=t+uX$,
ist schon von dem Erfinder auf eine scharfsinnige Art bewiesen worden\HistoricalNote{*)}{*)}{Hist.\ de l'Ac.\ de Berlin T.\ XXIV, Année 1768 (Berlin 1770), p.\ 275. Man findet die ganze Abhandlung übersetzt in dem dritten Bande, womit Hr.\ \textsc{Michelsen Eulers} Analysis des Unendlichen begleitet hat. Auch hat neuerlich Hr.\ D.\ \textsc{Murhard} diesen Beweis in einem eignen Programme durch einen gedrängten Auszug bekannter zu machen gesucht.}.
Allein da dieser Beweis sich darauf gründet, dass $X$ die Gestalt habe
\begin{center}
$\displaystyle ax^2+bx^3+cx^4+.\,.\,.\,.$\HistoricalNote{**)}{**)}{Eben das gilt von der Umkehrungsformel, die Hr.\ \textsc{Eschenbach} für sich erfunden und Hr.\ M.\ \textsc{Rothe} bewiesen hat, und die sich leicht mit dem \textsc{Lagrange}schen Satze verbinden lässt.}
\end{center}
\noindent und zugleich etwas weitläuftig ist, so haben verschiedene Geometer gesucht,
den Satz bloss aus der Natur der Functionen herzuleiten\HistoricalNote{***)}{***)}{\textsc{Lexell}, Comm.\ nov.\ Petr., T.\ XVI (1772), p.\ 230.\par
\textsc{Condorcet}, Miscell.\ Taurin., T.\ V (1770--73), Classe Math.\ p.\ 7.\par
\textsc{Cousin}, Astronomie physique, p.\ 15 (Paris 1787).\par
\textsc{Pfaff}, Archiv der r.\ u.\ a.\ M., 1.\ Band (1795), S.\ 81.}.
Von allen diesen Bemühungen ist indess bloss der Beweis des Hrn.\ Prof.\ \textsc{Pfaff}
völlig streng; die übrigen sind im Grunde blosse Induction. Ich hoffe, es werde
manchem Leser nicht unangenehm sein, hier einen andern gleichfalls strengen
Beweis zu finden, der mit dem des Hrn.\ Prof.\ \textsc{Pfaff} nichts gemein hat. Die
Hauptidee ist eigentlich wie in dem \textsc{Cousin}schen, aber die Ausführung scheint
mir neu und einfach. Verschiedenheit der Methoden dient immer, mehr Licht über
eine Wahrheit zu verbreiten, und gegenwärtiger Satz scheint unter uns bisher
weniger bekannt zu sein, als er es verdient.

